\documentclass[10pt,a4paper]{amsart}
\usepackage[margin=19mm]{geometry}
\usepackage{amsmath,amssymb,mathtools}
\usepackage[scr=rsfs,cal=euler]{mathalfa}
\usepackage{xfrac}
\usepackage[unicode,hidelinks,bookmarks=false]{hyperref}
\newcommand{\ie}{i.e.\ }
\newcommand{\cf}{cf.\ }
\newcommand{\resp}{respectively\ }
\newcommand{\Subsection}{\subsection}
\providecommand{\nobreakdash}{\nobreak\hbox{-}}
\setlength{\emergencystretch}{2em}
\multlinegap0pt
\usepackage{bm}
\usepackage{bbm}
\usepackage{dsfont}
\usepackage{xspace}
\usepackage{cancel}
\usepackage{mathtools}
\usepackage{amsthm}
\makeatletter
\@ifundefined{theorem}{\newtheorem{theorem}{Theorem}}{}
\makeatother
\newcommand{\dd}{\mathrm{d}}  % without the space
\newcommand{\va}{\mathbf{a}}
\newcommand{\vv}{{\mathbf{v}}}
\newcommand{\vx}{{\mathbf{x}}}
\newcommand{\vy}{{\mathbf{y}}}
\title[Microlocal analysis of a non-linear cone transform and appli]{Microlocal analysis of a non-linear cone transform and applications to Compton camera imaging}
\author{James W. Webber, Sean Holman}
\date{Source: arXiv:2608.03909v1 (CC BY 4.0)}
\begin{document}
\maketitle

\noindent\emph{Source and licence.} Microlocal analysis of a non-linear cone transform and applications to Compton camera imaging. Version arXiv:2608.03909v1 (CC BY 4.0), distributed under the Creative Commons Attribution 4.0 International licence, as recorded in the arXiv licence metadata for this exact version and shown on the abstract page https://arxiv.org/abs/2608.03909v1. Full rights notice: licenses/arxiv-2608.03909-CC-BY-4.0.txt.\par\smallskip

\noindent\textit{Editorial scope.} Counted unit: the Theorem whose source label is \texttt{natt} in the article (source0.tex lines 1209--1216, source label \texttt{natt}), delivered together with the article's own complete original proof of it (source0.tex lines 1217--1290, one \texttt{proof} environment of 74 lines). No mathematics is written, repaired or re-derived: statement and proof are the article's own text line for line. Required context retained uncounted: eq:Dkdef2 (lines 1193--1198). Every definition, display and label of those lines is retained unchanged. Omitted from this package and not redistributed: 1--1192, 1199--1208, 1291--1549. Nothing inside a retained excerpt is omitted or altered: every retained body is delivered exactly as the article's lines are written, so no omission receipt of any kind accompanies this package. Citations: the retained excerpts carry no citation command at all, so no bibliography entry of the article is transcribed and no reference is redistributed. Reference adaptation: none was needed and none was made; every reference inside the delivered unit resolves inside it, so no unresolved reference or citation is produced. Printed slips kept literal: no printed word, symbol, hypothesis or number of the counted statement or of its proof was altered. Layout and class: the article's own class is \texttt{amsart} with packages including lipsum, amssymb, amsthm, amsaddr, amsmath, tikz, graphicx, subcaption, enumitem, float, hyperref, placeins, pgfplots, fullpage; the wrapper is the lane's pinned \texttt{amsart} editorial wrapper at 10pt on A4, which restates the theorem-like environments the retained text needs and adds the article's own macro definitions that the retained text needs (\texttt{\textbackslash dd}, \texttt{\textbackslash va}, \texttt{\textbackslash vv}, \texttt{\textbackslash vx}, \texttt{\textbackslash vy}), with extra wrapper packages (bm, bbm, dsfont, xspace, cancel, mathtools). The delivered numbering is the unit's own. Assets: no retained span contains a figure environment, a graphics-inclusion command, a PDF-page inclusion, a table or a TikZ picture; no image file is copied into this package, the package declares no asset dependency and no image file is redistributed with it. Licence: version arXiv:2608.03909v1 of this article is distributed under the Creative Commons Attribution 4.0 International licence, as recorded in the arXiv licence metadata for this exact version and shown on the abstract page https://arxiv.org/abs/2608.03909v1. A full rights notice is delivered at licenses/arxiv-2608.03909-CC-BY-4.0.txt. Characters: every retained excerpt is pure ASCII, so the delivered engine, XeLaTeX with its default Latin Modern text font, reports no missing character.
\begin{equation}\label{eq:Dkdef2}
\begin{split}
\langle \widetilde{D}_k \mu, h \rangle_{L^2(I \times \Lambda_c)} & = \int_{\Lambda_c} \int_I \int_0^\infty \rho^{k} \mu(\va(s) + t \vx) h(s,\vx) \ \mathrm{d}s \mathrm{d}t \mathrm{d}\vx \\
& = \int_{\mathbb{R}^n} \mu(\vy)\left ( \int_I \int_0^\infty t^{k-n} h\left ( s, \frac{\vy-\va(s)}{t} \right )\ \mathrm{d}t \mathrm{d}s\right ) \mathrm{d}\vy
\end{split}
\end{equation}

\begin{theorem}
\label{natt}
Let $\Lambda$ be an open set on $S^{n-1}$, and let $A$ be a continuously
differentiable curve in $\mathbb{R}^n$ parametrized by $I \ni s \mapsto \va(s) \in A$. Suppose that the operator $\widetilde{D}_0: \mathcal{E}'(\mathbb{R}^n)\rightarrow \mathcal{D}'(I \times \Lambda_c)$ is defined as above using $A$ and $\Lambda$. Let $U \subset \mathbb{R}^n$ be bounded, open and disjoint from $A$. Assume that for each $\omega \in \Lambda$
there is an open interval $I_\omega$ such that for all $s \in I_w$, the ray $\{\va(s) + t\omega \ :\ t\geq 0\}$ misses $U$. If $\mu \in \mathcal{E}'(U)$ and
$\widetilde{D}_0\mu=0$, then $\mu = 0$ when restricted to the ``measured region"
$\{\vv+t\omega : \vv\in A, \omega \in \Lambda\}$.
\end{theorem}

\begin{proof}
% Let $A$ be a smooth embedded curve segment in $\mathbb{R}^2$ parametrised by $I \ni s \mapsto a(s)$, where $I$ is an open interval, and $S \subset S^1$ an open and connected set. Consider first the divergent beam moment transforms, which are defined for $f \in \mathbb{C}_c^\infty(\mathbb{R}^2)$ and $k\in \mathbb{Z}_{\geq 0}$ as
% \[
% \mathcal{D}_kf(s,\theta) = \int_0^\infty t^k f(a(s) + t \theta) \ \mathrm{d} t.
% \]
% \jc{\tred{update notation to be consistent with rest of paper.}}
% Let
% \[
% S_c =\left \{ x \in \mathbb{R}^2 \ : \ \frac{x}{|x|} \in S \right \}.
% \]
% Clearly $\mathcal{D}_k: C_c^\infty(\mathbb{R}^2) \rightarrow C^\infty(I \times S_c)$ continuously.

% First let us consider the definConsider a test function $h \in C_c^\infty(I\times S_c)$. Then
% \[
% \begin{split}
% \langle D_k f, h \rangle_{L^2(I \times S_c)} & = \int_{S_c} \int_I \int_0^\infty t^k f(a(s) + tx) h(s,x) \ \mathrm{d}x \mathrm{d}s \mathrm{d}t\\
% & = \int_{\mathbb{R}^2} f(y)\left ( \int_I \int_0^\infty t^{k-1} h\left ( s, \frac{y-a(s)}{t} \right )\ \mathrm{d}t \mathrm{d}s\right ) \mathrm{d}y
% \end{split}
% \]
% where we have changed variables in the second step and applied Fubini's theorem. From this calculation we see that formal adjoint $D_k^*$ is given by
% \[
% D_k^* h(y) =  \int_I \int_0^\infty t^{k-1} h\left ( s, \frac{y-a(s)}{t} \right )\ \mathrm{d}t \mathrm{d}s
% \]
% and it is clear $D_k^* : C_c^\infty(I \times S_c) \rightarrow C^\infty(\mathbb{R}^2)$ continuously. Based on this, we can extend $D_k$ to act on $f \in \mathcal{E}'(\mathbb{R}^2)$ by the formula
% \[
% \langle D_k f, h \rangle = \langle f, D_k^* h \rangle.
% \]
% The general theory says $D_k :  \mathcal{E}'(\mathbb{R}^2) \rightarrow \mathcal{D}'(I \times D_c)$ continuously.

Suppose $\mu$ satisfies the hypotheses of the theorem. We will prove by induction that $\widetilde{D}_k \mu = 0$ for all $k$, which we already know is true for $k= 0$. So, suppose $\widetilde{D}_{k-1} \mu = 0$ for some $k \geq 1$.

For $h(\cdot_s,\cdot_\vx) \in C_c^\infty(I \times S_c)$, using the identity
\[
\partial_s \left ( h \left ( s , \frac{\vy - \va(s)}{t} \right ) \right ) = \partial_{\widetilde{s}}|_{\widetilde{s} = s} h \left ( \widetilde{s} , \frac{\vy - \va(s)}{t} \right ) - \frac{\va'(s)^T}{t} \nabla_\vx h \left ( s , \frac{\vy - \va(s)}{t} \right )
\]
and a variation on integration-by-parts, we can calculate
\[
\begin{split}
\widetilde{D}_k^* [\partial_sh](\vy) & = \int_I \int_0^\infty t^{k-n} \partial_{\tilde{s}}|_{\tilde{s} = s} h\left ( \tilde{s}, \frac{\vy-\va(s)}{t} \right )\ \mathrm{d}t \mathrm{d}s\\
& =  \int_I \int_0^\infty t^{k-n-1} \va'(s)^T \nabla_\vx h\left ( s, \frac{\vy-\va(s)}{t} \right )\ \mathrm{d}t \mathrm{d}s \\
& = \widetilde{D}_{k-1}^* [(\va')^T \nabla_\vx h](\vy).
\end{split}
\]
Therefore,
\begin{equation}\label{eq:dDkfzero}
\langle \partial_s \widetilde{D}_k \mu, h \rangle = -\langle \mu,\widetilde{D}^*_k \partial_s h \rangle = -\langle \mu,\widetilde{D}^*_{k-1} (\va')^T\nabla_x h \rangle = -\langle \widetilde{D}_{k-1} \mu, (\va')^T\nabla_x h \rangle = 0.
\end{equation}
Now, by hypotheses we can find a covering of $\Lambda$ by open sets $\Lambda_i$ such that for each $i$ there is an open interval $I_i \subset I$ such that for all $\omega \in \Lambda_i$ and $s \in I_i$ the ray $\{\va(s) + t\omega \ :\ t\geq 0\}$ misses $U$. Let $\{\varphi_i\}$ be a partition of unity subordinate to the covering of $\Lambda$, and for $h \in C_c^\infty(I \times S_c)$ we have
\[
\langle \widetilde{D}_k\mu,h \rangle = \sum_{i} \langle \widetilde{D}_k\mu, \varphi_ih \rangle.
\]
We will prove each of the terms in the sum above is equal to zero. Now let $I_{s_0,\epsilon} = [s_0-\epsilon,s_0+\epsilon] \subset I_i$ and $\{s_j\}_{j=0}^N \subset I$ be a finite collection of points such that the intervals $I_{s_j,\epsilon} = [s_j - \epsilon,s_j+\epsilon]$ cover $I$. Let $\{\phi_j\}_{j=0}^N$ be a subordinate partition of unity on $I$. Then it is also sufficient to prove the terms
\[
\langle \widetilde{D}_k f,\varphi_i \phi_j h \rangle
\]
are all zero. For convenience, let us write \[
h_{ij} = \varphi_i \phi_j h
\]
and note that the support of $h_{ij}$ is contained in $I_{s_j,\epsilon} \times \Lambda_{i,c}$. Using the fundamental theorem of calculus, we have
\[
\begin{split}
\langle \widetilde{D}_k \mu, h_{ij} \rangle & = \left \langle \widetilde{D}_k \mu, -\int_0^{s_0-s_j} \partial_s h_{ij}(s-c,\vx) \ \mathrm{d}c + h_{ij}(s+s_j - s_0,\vx) \right \rangle\\
& = \int_{0}^{s_0-s_j} \langle \partial_s \widetilde{D}_k \mu, h_{ij}(s-c,\vx)\rangle \ \dd c + \langle \widetilde{D}_k \mu, h_{ij} (s+s_j-s_0,\vx) \rangle.
\end{split}
\]
The first term above is zero because of \eqref{eq:dDkfzero} while we can see the second is zero by considering \eqref{eq:Dkdef2} and noting that the support of $(s,\vx) \mapsto h_{ij}(s+s_j-s_0,\vx)$ is contained in $I_{s_0,\epsilon}\times \Lambda_{i,c} \subset I_i \times \Lambda_{i,c}$. This proves by induction that $\widetilde{D}_k \mu = 0$ for all $k \geq 0$.

By taking a partition of unity we can split up the measured region and so it is sufficient to prove
\[
\langle \mu, g \rangle = 0
\]
whenever the support of $g$ is contained in a small sliver in the measured region given by $\{ \va(s_0) + t \omega \ : \ \omega \in \Lambda_\epsilon\}$ where $\Lambda_\epsilon \subset \Lambda$ is sufficiently small. Approximating $g$ by polynomials in the radial variable centered at $a(s_0)$, which is possible by the Stone-Weierstrass theorem since the support of $\mu$ is compact, we can use the fact $\widetilde{D}_k 
\mu = 0$ for all $k$ to show that $\mu = 0$ when restricted to the measured region.
\end{proof}

\end{document}
